\documentclass[a4paper,12pt]{report}

% Basic packages for LuaLaTeX with Unicode support
% No need for inputenc with LuaLaTeX as it supports UTF-8 natively
\usepackage{fontspec}
\usepackage[brazil]{babel}
\usepackage{graphicx}
\usepackage{hyperref}
\usepackage{url}
\usepackage{amsmath}
\usepackage{amssymb}
\usepackage{listings}  % Regular listings is fine with LuaLaTeX
\usepackage{xcolor}
\usepackage{geometry}
\usepackage{indentfirst}
\usepackage{setspace}
\usepackage{cleveref}  % For enhanced cross-referencing
\usepackage{caption}   % For better caption handling
\usepackage{float}     % For improved figure/table placement
% \usepackage{luatextra} % Extra features for LuaLaTeX

% Hyperref configuration for better PDF output
\hypersetup{
    colorlinks=true,
    linkcolor=blue,
    filecolor=magenta,
    urlcolor=cyan,
    citecolor=green,
    pdftitle={Estudo acerca de memória no x86},
    pdfauthor={José Henrique Targino Dias Gois},
    pdfsubject={Memória no x86},
    pdfkeywords={sistemas operacionais, sistemas de memória, intel x86}
}

% Code listing styling
\definecolor{codegreen}{rgb}{0,0.6,0}
\definecolor{codegray}{rgb}{0.5,0.5,0.5}
\definecolor{codepurple}{rgb}{0.58,0,0.82}
\definecolor{backcolour}{rgb}{0.95,0.95,0.92}

\lstdefinestyle{mystyle}{
    backgroundcolor=\color{backcolour},
    commentstyle=\color{codegreen},
    keywordstyle=\color{magenta},
    numberstyle=\tiny\color{codegray},
    stringstyle=\color{codepurple},
    basicstyle=\ttfamily\footnotesize,
    breakatwhitespace=false,
    breaklines=true,
    captionpos=b,
    keepspaces=true,
    numbers=left,
    numbersep=5pt,
    showspaces=false,
    showstringspaces=false,
    showtabs=false,
    tabsize=2,
    extendedchars=true,
    texcl=true,
    inputencoding=utf8
}

\lstset{style=mystyle}

\geometry{margin=2.5cm}
\onehalfspacing

% Custom command for easier section referencing
\newcommand{\secref}[1]{Seção~\ref{#1}}
\newcommand{\figref}[1]{Figura~\ref{#1}}
\newcommand{\tabref}[1]{Tabela~\ref{#1}}
\newcommand{\eqnref}[1]{Equação~\ref{#1}}
\newcommand{\chapref}[1]{Capítulo~\ref{#1}}

% Document start
\begin{document}

% Cover page
\include{a_capa/a01_capa}

% Abstracts
\begin{abstract}
Este trabalho apresenta a implementação de um sistema de detecção de deadlocks e violação da ordem de aquisição de travas, inspirado no lockdep do kernel Linux. O sistema monitora operações de mutex em programas multithreaded, construindo um grafo de dependências entre travas e detectando potenciais situações de deadlock antes que ocorram. Utilizamos duas abordagens complementares: (1) detecção de ciclos no grafo de dependências, que indica um deadlock potencial, e (2) verificação da consistência na ordem de aquisição de travas entre diferentes threads. O sistema foi implementado como uma biblioteca compartilhada que utiliza a técnica de interposição de funções para interceptar chamadas à API pthread sem modificar o código fonte das aplicações monitoradas. Os testes realizados demonstram a eficácia do sistema na detecção de deadlocks clássicos como o problema AB-BA, onde duas threads tentam adquirir os mesmos mutexes em ordens diferentes. Este trabalho contribui para o desenvolvimento de ferramentas que aumentam a robustez de sistemas concorrentes, ajudando programadores a identificar problemas de sincronização difíceis de detectar.

\vspace{0.5cm}

\noindent\textbf{Palavras-chave:} deadlock, mutex, concorrência, detecção de erros, grafo de dependência, sistemas operacionais.
\end{abstract}

\begin{abstract}
\selectlanguage{american}
This work presents the implementation of a system for detecting deadlocks and lock acquisition order violations, inspired by the Linux kernel's lockdep. The system monitors mutex operations in multithreaded programs, building a dependency graph between locks and detecting potential deadlock situations before they occur. We use two complementary approaches: (1) cycle detection in the dependency graph, which indicates a potential deadlock, and (2) verification of consistency in the order of lock acquisition across different threads. The system was implemented as a shared library that uses function interposition to intercept calls to the pthread API without modifying the source code of monitored applications. Tests performed demonstrate the effectiveness of the system in detecting classic deadlocks such as the AB-BA problem, where two threads try to acquire the same mutexes in different orders. This work contributes to the development of tools that increase the robustness of concurrent systems, helping programmers identify synchronization problems that are difficult to detect.

\vspace{0.5cm}

\noindent\textbf{Keywords:} deadlock, mutex, concurrency, error detection, dependency graph, operating systems.
\selectlanguage{brazil}
\end{abstract}


% Custom TOC as defined in a04_toc.tex
\include{a_capa/a04_toc}

% Introduction
\chapter{Introdução}\label{chap:intro}

Este relatório apresenta o desenvolvimento de um sistema para detecção de deadlocks e violação da ordem de aquisição de travas em programas multithreaded. O sistema implementado, inspirado no lockdep do kernel Linux, monitora a aquisição de mutexes e verifica potenciais situações de deadlock antes que elas ocorram.

\section{Contextualização}

Em programação concorrente, o uso de mutexes (mutual exclusion) é essencial para garantir que apenas uma thread tenha acesso a recursos compartilhados por vez. No entanto, quando múltiplas threads tentam adquirir múltiplos mutexes em ordens diferentes, podem ocorrer deadlocks - uma situação onde duas ou mais threads ficam permanentemente bloqueadas, cada uma esperando por um recurso que outra possui.

Um exemplo clássico de deadlock ocorre quando:
\begin{itemize}
    \item Thread A adquire mutex1 e tenta adquirir mutex2
    \item Thread B adquire mutex2 e tenta adquirir mutex1
\end{itemize}

Este padrão, conhecido como AB-BA, cria uma dependência circular entre as threads, resultando em um deadlock. Em sistemas reais, estes deadlocks podem ser difíceis de reproduzir e diagnosticar devido à sua natureza dependente de timing.

\section{Objetivo}

O objetivo deste projeto é desenvolver um mecanismo de detecção de deadlocks que:

\begin{enumerate}
    \item Monitore a aquisição e liberação de mutexes
    \item Construa um grafo de dependência entre travas
    \item Detecte potenciais deadlocks verificando ciclos no grafo de dependência
    \item Alerte sobre violações na ordem de aquisição de travas
\end{enumerate}

\section{Abordagem}

Nossa implementação adota duas abordagens complementares:

\begin{enumerate}
    \item \textbf{Detecção de ciclos em grafo de dependências}: Construímos um grafo direcionado onde os nós são mutexes e as arestas representam a ordem de aquisição. Um ciclo neste grafo indica uma potencial situação de deadlock.

    \item \textbf{Verificação de ordem de aquisição}: Para cada thread, mantemos um histórico das ordens em que os mutexes foram adquiridos. Violações dessa ordem em aquisições futuras são sinalizadas como potenciais deadlocks.
\end{enumerate}

Para interceptar chamadas de mutex sem modificar o código fonte das aplicações, utilizamos a técnica de interposição de funções via LD\_PRELOAD, substituindo as chamadas padrão do pthread por nossas implementações instrumentadas.

% Implementation details
\chapter{Implementação}\label{chap:impl}

\section{Arquitetura do Sistema}

O sistema de detecção de deadlocks é composto por três componentes principais:

\begin{enumerate}
    \item \textbf{Núcleo de Detecção (lockdep\_core)}: Responsável pela lógica de detecção de deadlocks, mantendo o grafo de dependências e realizando a verificação de ciclos.

    \item \textbf{Interposição de Funções (pthread\_interpose)}: Intercepta chamadas às funções de mutex da biblioteca pthread e as redireciona para o núcleo de detecção.

    \item \textbf{Interface de API (lockdep.h)}: Define as estruturas de dados e a interface pública do sistema.
\end{enumerate}

\section{Estruturas de Dados}

Para representar o grafo de dependências e rastrear o estado das threads, utilizamos as seguintes estruturas:

\begin{itemize}
    \item \textbf{lock\_node\_t}: Representa um mutex no grafo de dependências.
    \item \textbf{dependency\_edge\_t}: Representa uma aresta direcionada no grafo (ordem de aquisição).
    \item \textbf{held\_lock\_t}: Mantém uma pilha dos mutexes atualmente possuídos por uma thread.
    \item \textbf{thread\_context\_t}: Armazena o contexto de uma thread, incluindo seus mutexes atualmente adquiridos.
\end{itemize}

\section{Algoritmos}

\subsection{Detecção de Ciclos}

Para detectar ciclos no grafo de dependências, implementamos um algoritmo de busca em profundidade (DFS):

\begin{enumerate}
    \item Para cada nó no grafo, iniciamos uma busca em profundidade.
    \item Durante a busca, marcamos os nós visitados para evitar processamento repetido.
    \item Se durante a busca encontramos um caminho que nos leva de volta ao nó inicial, temos um ciclo.
    \item A presença de um ciclo indica uma potencial situação de deadlock.
\end{enumerate}

\subsection{Verificação de Ordem de Aquisição}

Para cada thread, mantemos um histórico da ordem de aquisição de mutexes:

\begin{enumerate}
    \item Quando uma thread adquire um novo mutex, verificamos se este cria uma nova dependência de ordem.
    \item Se a nova dependência contradiz uma dependência existente, sinalizamos uma violação de ordem.
    \item Verificamos se a nova dependência cria um ciclo no grafo global de dependências.
\end{enumerate}

\section{Integração com Aplicações}

O sistema é projetado para ser transparente para as aplicações, não requerendo modificações no código fonte:

\begin{enumerate}
    \item Compilamos o sistema como uma biblioteca compartilhada (libpthread\_interpose.so).
    \item Utilizamos a variável de ambiente LD\_PRELOAD para carregar nossa biblioteca antes da biblioteca pthread padrão.
    \item As aplicações chamam as funções de mutex normalmente, mas nossas versões são executadas primeiro.
\end{enumerate}

\chapter{Resultados}\label{chap:results}

\section{Testes Realizados}

Para validar o sistema, desenvolvemos dois casos de teste principais:

\begin{enumerate}
    \item \textbf{simple\_test.c}: Demonstra a criação de dependências em uma única thread, adquirindo mutexes em uma ordem específica.
    \item \textbf{deadlock\_test.c}: Simula um cenário clássico de deadlock AB-BA, onde duas threads adquirem mutexes em ordens opostas.
\end{enumerate}

\section{Análise dos Resultados}

\subsection{Detecção de Violação de Ordem}

No teste de deadlock, o sistema detectou corretamente a violação de ordem quando:

\begin{enumerate}
    \item Thread 1 estabeleceu a ordem mutex1 → mutex2
    \item Thread 2 tentou estabelecer a ordem mutex2 → mutex1
\end{enumerate}

O sistema emitiu um alerta indicando que esta violação de ordem poderia levar a um deadlock.

\subsection{Detecção de Ciclo}

Após a violação de ordem, o sistema realizou a verificação de ciclo no grafo de dependências e detectou corretamente o ciclo formado pelas dependências:

\begin{enumerate}
    \item mutex1 → mutex2 (estabelecido pela Thread 1)
    \item mutex2 → mutex1 (tentado pela Thread 2)
\end{enumerate}

Este ciclo foi corretamente identificado como uma situação potencial de deadlock.

\section{Limitações}

O sistema atual possui algumas limitações:

\begin{enumerate}
    \item Não detecta deadlocks envolvendo outros tipos de primitivas de sincronização (semáforos, variáveis de condição).
    \item A sobrecarga de monitoramento pode ser significativa em aplicações com uso intensivo de mutexes.
    \item A interposição via LD\_PRELOAD não funciona com aplicações estaticamente vinculadas.
\end{enumerate}

% References
\chapter{Referências}\label{chap:refs}

\begin{thebibliography}{9}
\bibitem{lockdep}
  Rusty Russell,
  \emph{Lockdep: How to read its reports},
  Linux Kernel Documentation.

\bibitem{pthread}
  IEEE Computer Society,
  \emph{POSIX Thread Libraries},
  IEEE Std 1003.1-2017.

\bibitem{corbet}
  Jonathan Corbet,
  \emph{The kernel lock validator},
  LWN.net, 2006.

\bibitem{deadlock}
  Coffman, E.G., Elphick, M.J., Shoshani, A.,
  \emph{System Deadlocks},
  ACM Computing Surveys, 1971.
\end{thebibliography}
% \include{x_referencias/x01_referencias}

% Appendices
\appendix
\include{j_apendices/j01_appendix}

\end{document}
